% !TEX root = ../main.tex
\section{Claim-to-Evidence Map}
\label{ch:appendix-claims}

This appendix maps each hypothesis, research question, and empirical claim to the relevant evidence in Chapter~\ref{ch:results}.

\dissertationSubheading[sec:map-h1]{H1 / RQ1: Rank divergence}

\begin{itemize}
\item Statement: on the same cohort, EndorseRank and AWP rank wallets differently.
\item Evidence: Section~\ref{sec:rank-divergence}; Table~\ref{tab:tie-sensitivity}, whose last row gives the inter-method $\tau_b$ with its interval, also with the wallets outside each graph kept as isolated nodes; Figures~\ref{fig:rank-scatter} and~\ref{fig:score-distributions}.
\item Interpretation: the rankings are not equivalent, and their agreement comes from the wallets that EndorseRank scores zero; the O1 criterion is met (Table~\ref{tab:objectives-achievement}).
\end{itemize}

\dissertationSubheading[sec:map-h2]{H2 / RQ2: Same-window alignment}

\begin{itemize}
\item Statement: in the same window, each score agrees with the checks built from the edges it walks on.
\item Evidence: the family means of every score with their CIs (Table~\ref{tab:alignment-hybrid}); the same-window coefficients with the wallets outside a graph kept as isolated nodes (Table~\ref{tab:tie-sensitivity}).
\item Note: the agreement is partly built in, since the proxies come from the same events as the graphs. EndorseRank's allowance and transfer coefficients have different ceilings (Section~\ref{sec:same-window-alignment}), so the claim says nothing about which of the two it follows more closely, and the isolated-node rule was tried after the results were known. None of these comparisons was fixed in advance. The paired differences between EndorseRank and AWP and their per-proxy results are in Appendix~\ref{ch:appendix-er-awp}.
\end{itemize}

\dissertationSubheading[sec:map-h3]{H3 / RQ3: Computational cost}

\begin{itemize}
\item Statement: at the same $n$, a score's computation time follows the number of edges in its graph, so EndorseRank, on the sparsest graph, runs fastest.
\item Evidence: Table~\ref{tab:benchmark-runtime} (runtime, SD, memory, iterations, $|V|$, $|E|$); Tables~\ref{tab:benchmark-scaling-incohort} and~\ref{tab:benchmark-scaling}; Figure~\ref{fig:runtime-scaling}.
\item Note: the expanded-pool stages add sampled addresses but no edges, and hence no solver-graph nodes, beyond the matched-cohort graph (Section~\ref{sub:runtime-scaling}); the claim is confined to edge count.
\end{itemize}

\dissertationSubheading[sec:map-rq4]{RQ4: Temporal holdout}

\begin{itemize}
\item Statement: with scores fixed at $t_1$, the scores predict the new owner--spender approvals and new transfer senders that appear after $t_1$ and up to $t_2$ on the spender cohort ($n=1{,}335$), with the edge type deciding which label a score ranks well, but no PageRank predicts either label better than the raw $t_1$ degree of its own edge type.
\item Evidence: Tables~\ref{tab:holdout-spenders} and~\ref{tab:holdout-diff} (Section~\ref{sec:holdout-results}); the June--August window, Table~\ref{tab:fresh-posthoc}; the spenders that had received a transfer before the freeze, Table~\ref{tab:holdout-receiving}.
\item Qualification: both labels are the constructs as observed in a later period (new approvals, new senders), not credit or trading outcomes, and both are counts of arrivals. RQ4 replaced the proposal's trading-success validation (Section~\ref{sub:gf-pr}). Apart from the increment of C-PR ($\lambda=1$), none of these comparisons was fixed in advance. The differences between EndorseRank and AWP are in Appendix~\ref{ch:appendix-er-awp}.
\end{itemize}

\dissertationSubheading[sec:map-rq5]{RQ5: Coupled operator}

\begin{itemize}
\item Statement: under the rule of 14~September, a per-node convex mixture of the raw-amount allowance and transfer layers (C-PR) does not beat both layers walked alone, and the seeded ablation (S-PR) reproduces the transfer walk. The registered replication does not confirm that C-PR improves on the transfer walk alone, as it did in an exploratory comparison on the spring holdout: F1 and F3 hold, F2 fails. With AWP as the comparator, F2 fails by a wider margin.
\item Evidence, rule of 14~September: Table~\ref{tab:holdout-diff} for the primary part, and Tables~\ref{tab:alignment-hybrid} and~\ref{tab:tau-diff-hybrid} for the secondary part (Section~\ref{sec:coupled-results}).
\item Evidence, transfer walk alone: the exploratory spring contrast in Table~\ref{tab:holdout-diff}, and the registered decision (F1--F3, with F4--F6 reported beside it) in Table~\ref{tab:fresh-contrasts} (Section~\ref{sub:fresh-results}).
\item Evidence, AWP: the registered sensitivity analysis in Table~\ref{tab:fresh-contrasts}, and the spring contrasts computed after the labels were known in Table~\ref{tab:holdout-diff}.
\item Evidence, cost: runtime rows of Table~\ref{tab:benchmark-runtime}.
\item Qualification: tested with fixed $\lambda\in\{0.25,0.5,0.75\}$, raw-amount layers, uniform teleportation, transition-level mixing and one chain. The rule of 14~September and the values of $\lambda$ were set in the configuration before any coupled score was computed (Section~\ref{sec:holdout-protocol}); the rule for the comparison with the transfer walk alone was written on 28~September~2026 for the June--August window and registered on 1~October~2026 together with two sensitivity analyses, one with AWP as the comparator, before the extraction (Section~\ref{sub:fresh-holdout}). The registration documents call the two layers walked alone EndorseRank and AWP.
\end{itemize}

\dissertationSubheading[sec:map-c1c4]{Claims C1--C5}

\begin{itemize}
\item C1: the runtime evidence listed under H3/RQ3; the profiling run and the edge ratio are in Section~\ref{sec:benchmark-results} and Appendix~\ref{sec:pr-sparsity}. The claim is specific to EndorseRank, since C-PR costs at least as much as AWP.
\item C2: the same-window agreement of every score with the checks of its own edges (Table~\ref{tab:alignment-hybrid}), partly built in and, for EndorseRank, set by the rule for wallets outside the graph (Table~\ref{tab:tie-sensitivity}).
\item C3: the rank divergence listed under H1/RQ1 (Table~\ref{tab:tie-sensitivity}, Figure~\ref{fig:rank-scatter}).
\item C4: the comparison of every PageRank with the raw degree of its own edge type on the spender holdout (Tables~\ref{tab:holdout-spenders} and~\ref{tab:holdout-diff}), with the coverage check (Table~\ref{tab:holdout-receiving}), the exploratory trader run as its limit on the traders (Table~\ref{tab:holdout-traders}) and the June--August window as its repetition (Table~\ref{tab:fresh-posthoc}); the robustness checks and the sample-definition sensitivity are reported in Tables~\ref{tab:robustness-damping}--\ref{tab:robustness-sample-size}, Table~\ref{tab:endorserank-restarts} and Figure~\ref{fig:damping-sensitivity}.
\item C5: coupling does not beat its two layers walked alone, a bounded negative result, and the registered replication does not confirm its exploratory gain over the transfer walk alone (Table~\ref{tab:fresh-contrasts}); see RQ5 above and Section~\ref{sec:coupled-results}.
\end{itemize}

\dissertationSubheading[sec:map-nonclaims]{Explicit non-claims}

This research does \emph{not} claim:

\begin{itemize}
\item that any of the scores is the better measure of reputation in general, or that AWP has become obsolete (the comparisons between EndorseRank and AWP in Appendix~\ref{ch:appendix-er-awp} follow the constructs and depend on coverage);
\item that EndorseRank aligns more closely with the allowance family than with the transfer family (the two coefficients have different ceilings);
\item that any PageRank of this thesis forecasts the future growth of its own construct better than the raw $t_1$ degree does;
\item that allowance edges lead transfer edges on liquidation avoidance at the PageRank level, or that EndorseRank beats AWP there (the PageRank contrasts of Section~\ref{sub:neutral-label-results} are not resolved);
\item that any score measures creditworthiness, default risk or trading skill (the allowance signal is lower on the profitable share, and the data contain no default label);
\item that the liquidation-avoidance signal belongs to the allowance edge and not to account type (almost every approval-receiving trader is a contract account, and account type is unresolved);
\item that the value parameter $b=1$, or EndorseRank's uniform restart, is the best choice (both were fixed by this study, the restart after comparing the two rules on these data);
\item that C-PR improves on a walk over the transfers alone (the registered rule failed on new senders) or on AWP (it falls behind AWP on new senders in both windows);
\item that no combination of the two edge types can beat the single-layer methods, only that the per-node convex mixture tested here, with fixed $\lambda$, does not;
\item that runtime scales favorably when the edge population grows beyond the cohort graph;
\item that, without replication, the results generalize beyond Arbitrum One and the study window;
\item that on-chain scores substitute for collateral or legal identity systems.
\end{itemize}

\dissertationSubheading[sec:map-artifacts]{Replication artifacts}

\begin{enumerate}
\item A merged wallet-ranking table with the scores of EndorseRank, AWP, C-PR ($\lambda\in\{0,0.25,0.5,0.75,1\}$) and S-PR.
\item The machine-readable holdout summaries for the spender cohort and the matched traders, with the paired contrasts and primary-criterion verdict (committed under \texttt{data/3-published-results-for-thesis/spring-holdout-2026-03-to-2026-05/}).
\item The machine-readable evaluation summary: authoritative same-window numbers, per-run timings and the bootstrap protocol (committed under \texttt{data/3-published-results-for-thesis/}).
\item The registered replication: the registration file and plan as committed on 1~October~2026, the extraction manifest with the hashes of both files, the evaluation summary with every coefficient, contrast and verdict, and the post hoc scores of EndorseRank on the same cohorts (committed under \texttt{config-copies/} and \texttt{registered-replication-2026-06-to-2026-08/} in \texttt{data/3-published-results-for-thesis/}).
\item The exported table fragments (\texttt{results/tables/}).
\item The exported vector figures (\texttt{Figures/generated/}).
\item The extracted logs and the tables built from them (\texttt{data/1-raw-blockchain-logs/} and \texttt{data/2-processed-tables-and-evaluations/}); files over 100\,MiB are left out and can be rebuilt from the monthly files.
\item The archived extended-baseline matrices (\texttt{data/4-archived-extended-baselines/}, optional).
\end{enumerate}

The tables and figures are rebuilt by running the evaluation pipeline on real data with the robustness sweeps, exporting the tables, and then redrawing the figures from the same artifacts (Appendix~\ref{sec:pipeline-stages}).

\dissertationSubheading[sec:map-c6]{C6 --- Outcomes built from neither edge type}

\begin{itemize}
\item Statement: on the trader outcome built from neither edge type that bears on counterparty risk, later liquidation avoidance, the approval count at the freeze ranked traders above the transfer count in the registered window, after stratification by the number of closes, and in the spring window.
\item Evidence: Tables~\ref{tab:neutral-label-contrasts} and~\ref{tab:neutral-label-recipients} (Section~\ref{sub:neutral-label-results}); all six trader labels in both windows, Table~\ref{tab:neutral-label-grid}; the protocol, Section~\ref{sub:neutral-label-protocol}; the plan and claim ledger, \texttt{analysis/neutral-label/PLAN.md}, committed in \texttt{20e250f} before the analysis ran, with its script and results in the same folder.
\item Qualification: a pre-specified post hoc analysis of existing data that adds no research question. The claim is count-level, since the like-for-like PageRank pair points the same way but is not resolved. Most recipients are contract accounts, so account type is an unresolved confound, and the same signal is lower on the profitable share, so it ranks forced-liquidation avoidance, not trading success. Any use in screening would first need a registered test on new data with account type controlled (Section~\ref{sec:future-work}).
\end{itemize}
